% Generated draft: requires mathematical and source review.
\documentclass[11pt]{article}
\usepackage[margin=25mm]{geometry}
\usepackage{fontspec}
\setmainfont{Latin Modern Roman}
\usepackage{amsmath,amssymb,mathtools,unicode-math}
\setmathfont{Latin Modern Math}
\usepackage{enumitem,needspace,xurl,hyperref}
\setlength{\parindent}{0pt}
\setlength{\parskip}{5pt}
\setlength{\emergencystretch}{3em}
\newcommand{\N}{\mathbb{N}}
\newcommand{\Z}{\mathbb{Z}}
\newcommand{\Q}{\mathbb{Q}}
\newcommand{\R}{\mathbb{R}}
\newcommand{\C}{\mathbb{C}}
\newcommand{\F}{\mathbb{F}}
\newcommand{\A}{\mathbb{A}}
\newcommand{\PP}{\mathbb{P}}
\newcommand{\E}{\mathbb{E}}
\newcommand{\eps}{\varepsilon}
\newcommand{\dd}{\,\mathrm{d}}
\begin{document}
\section*{UnsolvedMath ID: 30003847}
\textbf{Source problem code:} catalogue entry
\subsection{Definitions and mathematical statement}
Let \(H\) denote the ambient group written \(\mathrm{PL}_{o}(I)\) in the supplied record, and put
\[
X=\{G\leq H: G\text{ is finitely generated}\}.
\]
For groups \(G,K\), write \(G\hookrightarrow K\) if there exists an injective group homomorphism from \(G\) to \(K\). Define biembeddability by
\[
G\sim K\quad\Longleftrightarrow\quad G\hookrightarrow K\ \text{and}\ K\hookrightarrow G.
\]
Write \([G]\) for the equivalence class of \(G\) and define the partially ordered set
\[
[X]=X/{\sim},\qquad [G]\preceq[K]\quad\Longleftrightarrow\quad G\hookrightarrow K.
\]
Here embeddings are read as abstract group embeddings; the record does not specify any additional compatibility with the actions on \(I\).

The supplied extraction combines two questions:
\begin{enumerate}
\item Can the structure of the biembeddability poset \(([X],\preceq)\) be understood further?
\item Does Cleary's golden-ratio Thompson group \(F_{\tau}\) embed in Thompson's group \(F\)? In symbols, does there exist an injective homomorphism
\[
\varphi:F_{\tau}\longrightarrow F?
\]
\end{enumerate}
The first question specifies no particular structural property to establish and therefore does not determine a precise conjecture. The record does not define \(I\), the notation \(\mathrm{PL}_{o}(I)\), or the groups \(F\) and \(F_{\tau}\) by presentations or actions. The poset construction above is consequently relative to the intended ambient group, and the second question retains the named groups without supplying missing definitions.

The supplied statement-verification note reports that these are independent numbered questions merged during extraction, and marks the combined statement as \emph{unrecoverable}. The supplied dated literature assessment reports a negative answer to the embedding question by Hyde--Moore, while retaining an \emph{open} label for the combined record because the structural question is incomplete. These are supplied curation claims, not independently verified mathematical or literature conclusions.
\subsection{Short English statement}
The record combines a broad question about the ordering of finitely generated subgroups by mutual embeddability with the question whether Cleary's golden-ratio Thompson group embeds in Thompson's group \(F\). The structural target and several definitions are missing. Supplied dated curation reports that the embedding question has a negative answer, but keeps the combined record labeled open; this draft requires human mathematical and source review.
\subsection{Sources}
\begin{itemize}
\item Local UnsolvedMath registry record, permanent ID \(30003847\), title \emph{Biembeddability and the Golden-Ratio Thompson Group}, collection \emph{Oberwolfach Reports Open Problems}, exact source code \texttt{catalogue entry}. Repository registry locations are \texttt{data/unsolvedmath/problems.json} and \texttt{catalog/problems.json}; this draft uses the supplied complete record and does not assert which file contains it. The record gives proposed year 2018 and names no proposer. It attributes the golden-ratio Thompson group to Cleary.
\item Supplied report citation: \emph{Cohomological and Metric Properties of Groups of Homeomorphisms of R} (2019), Oberwolfach Reports, DOI 10.4171/owr/2018/26, pp. 1579--1633. \url{https://doi.org/10.4171/owr/2018/26}. The background describes Oberwolfach Reports as documenting mathematical workshops and their open-problem discussions.
\item Supplied literature source: Hyde--Moore, \emph{Subgroups of PL+ I}. \url{https://arxiv.org/abs/2103.14911}.
\item Supplied literature source: \emph{Oberwolfach Report 26/2018}. \url{https://ems.press/content/serial-article-files/46748?nt=1}.
\item Literature assessment dated exactly \texttt{2026-08-21}: supplied dated curation reports that Hyde--Moore answer the \(F_{\tau}\)-into-\(F\) embedding question negatively, while an incomplete structural question leaves the combined record labeled open. The background qualifies this assessment as dated literature triage derived from supplied OWR curation and requests independent verification of time-sensitive status and cited claims. Neither that assessment nor the supplied claim that the official report contains two independent numbered questions has been independently verified here.
\end{itemize}
\end{document}
